% SPDX-License-Identifier: CC-BY-SA-4.0
% SPDX-FileCopyrightText: Louis Royer <infos.louis.royer@gmail.com>
\documentclass[../main.tex]{subfiles}
\begin{document}
\begin{licenseHeader}
L’ouverture des connaissances scientifiques est l’un des piliers de la science ouverte, et le libre accès aux publications scientifiques est nécessaire à cette ouverture.

Le 2\textsuperscript{ème} plan national pour la science ouverte, établi en 2021, a pour ambition de généraliser les pratiques de science ouverte en France, et la loi de programmation de la recherche fixe l’objectif de 100~\% de publications en accès ouvert en 2030.

\vspace{0.5cm}
\astersmall
\vspace{0.5cm}

Pour toutes ces raisons, et par ferme conviction que la généralisation de la science ouverte constitue l’avenir de la recherche, cette thèse est publiée sous une licence libre.
\end{licenseHeader}

\vspace{2cm}
\asterthree
\vspace{2cm}
L’ensemble des fichiers source \LaTeX{} de ce manuscrit, ainsi que les données brutes des mesures effectuées sont disponibles à l’adresse: \url{https://github.com/louisroyer/these}

\begin{licenseFooter}
\faCreativeCommons \hspace{0.1cm} \faCreativeCommonsBy \hspace{0.1cm} \faCreativeCommonsSa \hspace{0.5cm} Ce document est mis à disposition selon les termes de la \href{https://creativecommons.org/licenses/by-sa/4.0/deed.fr}{licence CC BY-SA 4.0}. \\
\url{https://creativecommons.org/licenses/by-sa/4.0/}
\end{licenseFooter}
\end{document}